\subsection{Einsatz einer Zenerdiode zur Spannungsstabilisierung}
% Alt: Aufgabe 6
\Aufgabe{
  \label{Aufg6}
  Mit Hilfe einer Zenerdiode soll ein Lastwiderstand \( R_\mathrm{L} = \mathrm{160\,\Omega} \) mit einer stabilisierten Spannung von \( U_\mathrm{A} = \mathrm{8\,V} \) versorgt werden. \\
  Dem Datenblatt der Zenderdiode ist zu entnehmen, dass bei einer Zenerspannung von \( \mathrm{-8\,V} \) ein Zenerstrom von \( \mathrm{-0{,}45\,A} \) fließt.

  \begin{figure}[H]
    \centering
    \input{Tikz/src/FigZDiodeStabilisierung}\alttext{Zu sehen ist eine elektrische Schaltung mit vier Bauteilen zur Bestimmung Ermittlung der Spannungsstabilisierung einer Zenerdiode. Auf der linken Seite befindet sich die Spannungsquelle \(U_\mathrm{E}\), deren Spannungspfeil nach unten zeigt. In Reihe zur Spannungsquelle sind ein Widerstand \(R\) und eine Zenerdiode \(D_\mathrm{z}\) geschaltet.  Parallel zu der Zenerdiode befindet sich ein Lastwiderstand gekennzeichnet mit \(R_\mathrm{L}\). Ein Spannungspfeil in Stromrichtung kennzeichnet die Spannung der Parralelschaltung mit Der Beschriftung \(U_\mathrm{A}\).}
    \label{fig:FigZDiodeStabilisierung}
  \end{figure}

  \begin{itemize}
    \item[a)] Berechnen Sie den Vorwiderstand \( R_\mathrm{V} \), sodass sich bei einer Eingangsspannung von \( U_\mathrm{E} = \mathrm{10\,V} \) eine Ausgangsspannung von \( U_\mathrm{A} = \mathrm{8\,V} \) einstellt.
    \item[b)] Wie groß sind die Leistungen, die an \( R_\mathrm{V} \), \( R_\mathrm{L} \) und an der Diode \( D_\mathrm{Z} \) umgesetzt werden?
  \end{itemize}
}


\Loesung{
  \begin{figure}[H]
    \centering
    \input{Tikz/src/LsgZDiodeStabilisierung}
    \alttext{Die Abbildung zeigt eine Schaltung mit einer Eingangsspannungsquelle \(U_\mathrm{E}\), einem Vorwiderstand \(R_\mathrm{V}\), einer Z-Diode \(D_\mathrm{D}\), sowie einem Laswtwiderstand \(R_\mathrm{L}\).
    Links befindet sich die Spannungsquelle \(U_\mathrm{E}\), die zwischen dem oberen linken Knoten und Masse geschaltet ist. Vom oberen Anschluss der Spannungsquelle führt der Widerstand \(R_\mathrm{V}\) zum rechten oberen Knoten. Der durch diesen Widerstand fließende Strom ist mit \(I\) gekennzeichnet. Die Z-Diode liegt parallel zur Ausgangsklemme.
    Zwischen dem oberen und dem unteren mittleren Knoten ist eine Z-Diode \(D_\mathrm{D}\) geschaltet. Der durch die Z-Diode fließende Strom ist als Zenerstrom \(I_\mathrm{Z}\) eingezeichnet.
    Rechts ist der Lastwiderstand \(R_\mathrm{L}\) zwischen dem oberen rechten Knoten und Masse angeschlossen. Über diesem Widerstand ist die Spannung \(U_{\mathrm{RL}}\) markiert.
    Die Ausgangsspannung \(U_\mathrm{A}\) wird zwischen dem oberen rechten Knoten und Masse gemessen.
    Die beiden eingezeichneten Umlaufpfeile kennzeichnen zwei Maschen:
    Links \(M_1\) umfasst die Spannungsquelle, den Vorwiderstand und die Z-Diode.
    rechts \(M_2\) umfasst den Vorwiderstand, die Z-Diode und den Lastwiderstand.}
    \label{fig:LsgZDiodeStabilisierung}
  \end{figure}
  \begin{itemize}
    \item[a)]
          \begin{gather*}
            I_\mathrm{RL} = \frac{\mathrm{8\,V}}{\mathrm{160\,\Omega}} = \mathrm{0{,}05\,A} \\
            I = \mathrm{0{,}05\,A} + \mathrm{0{,}45\,A} = \mathrm{0{,}5\,A} \\
            U_\mathrm{RV} = \mathrm{10\,V} - \mathrm{8\,V} = \mathrm{2\,V} \\
            R_\mathrm{V} = \frac{\mathrm{2\,V}}{\mathrm{0,5\,A}} = \mathrm{4\,\Omega}
          \end{gather*}

    \item[b)]
          \begin{gather*}
            P_\mathrm{RV} = \mathrm{2\,V} \cdot \mathrm{0{,}5\,A} = \mathrm{1\,W} \\
            P_\mathrm{RL} = \mathrm{8\,V} \cdot \mathrm{0{,}05\,A} = \mathrm{0{,}4\,W} \\
            P_\mathrm{DZ} = \mathrm{8\,V} \cdot \mathrm{0{,}45\,A} = \mathrm{3{,}6\,W}
          \end{gather*}
  \end{itemize}
  Siehe: Abschnitt \ref{sec:SpezielleDioden}
}